\section{Introduction}
\label{sec:intro}

Dynamic Bayesian networks (DBNs) describe multivariate time series by a directed acyclic
graph (DAG) of contemporaneous (intra-slice) effects together with lagged (inter-slice)
effects. In the linear case this is a structural vector autoregression,
\begin{equation}
  \label{eq:svar}
  x_t^\top = x_t^\top W + \sum_{k=1}^{p} x_{t-k}^\top A_k + z_t^\top, \qquad t = 1,\dots,n,
\end{equation}
where $x_t\in\R^d$, $W\in\R^{d\times d}$ is the weighted adjacency matrix of a DAG,
$A_1,\dots,A_p\in\R^{d\times d}$ are unconstrained and $z_t$ is noise with independent
coordinates \citep{pamfil2020dynotears}. Learning $(W, A)$ from data is NP-hard already in
the static case \citep{chickering2004large}. Two lines of work dominate. \emph{Exact}
methods based on mixed-integer programming (MIP) certify global optimality of a score and,
on small instances, recover the true structure essentially perfectly; \exdbn{}
\citep{rytir2025exdbn} is the state of the art for DBNs. \emph{Continuous} methods such as
DYNOTEARS \citep{pamfil2020dynotears}, built on NOTEARS \citep{zheng2018notears}, or
\psidag{} \citep{ziu2024psi} scale to many more variables but come without global guarantees.

In many applications the time series are not held by one party. Hospitals hold the records of
their own patients, banks hold their own exposures (the credit default swap application of
\citet{rytir2025exdbn} is of this kind), and pooling the data is ruled out by regulation or
competition. Federated learning \citep{mcmahan2017fedavg,kairouz2021advances} addresses this
by exchanging models instead of data. Federated structure learning has been studied for
static DAGs \citep{ng2022federated,gao2023feddag,li2024fedcdh}, but to our knowledge not for
DBNs, and the following question has not been quantified: \emph{how much statistical accuracy
does one give up by learning a DBN in a federated, first-order fashion, rather than by solving
the exact MIP on pooled data, and what does one gain in scalability?} We call the loss of
accuracy the \emph{price of federation}.

\paragraph{Contributions.}
\begin{enumerate}
  \item We propose \feddbn{} (\cref{sec:method}), an extension of \psidag{} \citep{ziu2024psi}
  to DBNs (an acyclic intra-slice matrix together with unconstrained lag matrices) and to the
  federated setting. Clients run local accelerated gradient steps in the spirit of FedAvg. The
  server averages the models and performs \psidag's order-based projection, which needs no
  data. Only parameter matrices cross the network.
  \item We observe that for the linear least-squares score the \exdbn{} MIP depends on the data
  only through the pooled Gram matrix, so it can itself be federated in a single round of
  secure aggregation (\cref{prop:gram}). The statistical comparison against \exdbn{} on pooled
  data is therefore a comparison against the best \emph{exact federated} method, and the
  reformulation removes the dependence of the MIP model size on $n$.
  \item We define the price of federation and decompose it into an estimator term, an
  \emph{ordering} term (the topological order chosen by the projection), an \emph{incomplete
  optimisation} term and a \emph{federation} term, together with the cost of
  \emph{communication}: rounds, latency, bandwidth and staleness (\cref{sec:price}).
  \item We explain why MIP does not scale, sequentially or in parallel
  (\cref{sec:mip}). Even solving only the root relaxation means solving an instance of
  a P-complete problem class \citep{dobkin1979lp,serna1991approximating}, which is unlikely to
  be efficiently parallelisable. The weak big-$M$ relaxation means the root node bound is
  close to useless. Parallel branch-and-bound has historically scaled poorly
  \citep{shinano2010parascip,koch2012million,bernalneira2024survey}.
  \item We measure the price of federation on \exdbn's own benchmark, with \exdbn's generator,
  parameter grid and evaluation code, and report scalability to $d=1000$ variables
  (\cref{sec:experiments}).
\end{enumerate}
